\section{Stress and rotations}\label{sec:stress}

We compare the fields of Section~\ref{sec:continuum} with Maxwell fields at small distances, then normalize a conserved tensor by its shell fluxes. The volume refinement and the table~$S$ remain fixed throughout. Subsequences taken below concern only comparison radii.

The comparison permits the basis of local fields to vary with the radius. This accommodates mixing among fields in a fixed finite-dimensional space.
\begin{definition}[Scaling frames]\label{def:D12}
A scaling sequence consists of $r_k>0$, $r_k\to0$, and isomorphisms
\begin{equation*}
\scaleframe{k}{d}:\maxwellspace{d}\longrightarrow F_d\,.
\end{equation*}
It is stable if, at every fixed degree horizon, these maps eventually preserve conjugation, $\hypercubic$, units and inclusions exactly. Scaling dimensions are included in the maps: an ordinary comparison insertion is $\scaleframe{k}{d}u$ at~$r_kx$, with no further factor~$r_k^d$.
\end{definition}

The first comparison concerns correlations integrated against joint tests. A single sequence works for all fixed lists of fields.
\begin{proposition}[Maxwell comparison]\label{prop:MC}
There is one stable scaling sequence such that, for every $n,u_i\in\maxwellspace{d_i}$, compact $\mathcal K\subset\Conf{n}$, $M\ge0$ and $\varepsilon>0$, some $k_0(n,u,\mathcal K,M,\varepsilon)$ gives, for $k\ge k_0$ and $f\in C_c^\infty$ supported in~$\mathcal K$ with $\|f\|_\infty\le M$,
\begin{equation*}
\left|\int f(x)\big[S_n(\scaleframe{k}{d_i}u_i;r_kx)-\maxwellcorr_n(u;x)\big]\,dx\right|\le\varepsilon M\,.
\end{equation*}
\end{proposition}
The proof is in Section~\ref{sec:partII-G}, II.G.2.1. The lattice comparison stops near the requested physical scale and compares the resulting correlations with the free Gaussian ones. Finitely many separating correlations make the comparison maps injective; the dimension equality of FC makes them isomorphisms. Finite-dimensional corrections give the exact eventual compatibilities in D12.

To read translation and angular charge, integrate a conserved tensor through a shell surrounding just one insertion. Let $A\in F_d$, let $\boldsymbol B$ be a spectator field list at~$x$, and take $(x_0,x)\in\Conf{n+1}$.
\begin{definition}[Shell fluxes]\label{def:D13}
For $J_{ab}\in F_4$ and $K_{\rm rot}\in\mathfrak{so}(4)$, choose real $\chi\in C_c^\infty$ equal to one near~$x_0$ and supported away from spectators. Define
\begin{align*}
\mathcal P_b^J&=-\sum_a\int\partial_a\chi(y)
S(J_{ab}(y),A(x_0),\boldsymbol B(x))\,dy,\\
\mathcal Y_{K_{\rm rot}}^J&=-\sum_{a,b}\int\partial_a\chi(y)(K_{\rm rot}(y-x_0))_b
S(J_{ab}(y),A(x_0),\boldsymbol B(x))\,dy\,.
\end{align*}
These fluxes avoid every insertion point.
\end{definition}

The centered real symmetric $\hypercubic$-tensorial conserved tensors of degree four form a real line, as proved in II.G.2.2. For a nonzero member~$J$, II.G.2.3 constructs linear charge maps $\mathsf p_b:F_d\to F_{d+1}$ and $\mathsf Y_{K_{\rm rot}}:F_d\to F_d$ by
\begin{equation*}
\mathcal P_b^J=S((\mathsf p_bA)(x_0),\boldsymbol B(x)),\qquad
\mathcal Y_{K_{\rm rot}}^J=S((\mathsf Y_{K_{\rm rot}}A)(x_0),\boldsymbol B(x))\,.
\end{equation*}
Here centered means that every component has zero one-point function. The charge maps are complex linear in~$A$, and the angular map is real linear in~$K_{\rm rot}$.

Section~\ref{sec:partII-G}, II.G.3.1, identifies these maps with differentiation and rotations. It constructs continuous complex linear representations $R_d$ of~$O(4)$, extending~$\hypercubic$, fixing units and commuting with conjugation and inclusions. For a real $c_\star\ne0$ common to every degree, the generators and charge maps satisfy
\begin{equation*}
\mathsf L_d(K_{\rm rot})=\left.\frac{d}{dt}R_d(e^{tK_{\rm rot}})\right|_{t=0},
\qquad \mathsf p_b=c_\star\partial_b,\qquad
\mathsf Y_{K_{\rm rot}}=c_\star\mathsf L_d(K_{\rm rot})\,.
\end{equation*}
The proof controls exterior angular flux before extending the rotation identity to all separated configurations.

This common constant fixes the normalization of the stress tensor by its action on fields.
\begin{theorem}[Stress tensor]\label{prop:ST}\label{thm:ST}
There is a real symmetric $O(4)$-tensor $T_{ab}\in F_4$ with $\sum_a\partial_aT_{ab}=0$ such that, for every degree, field, spectator list, separated configuration and cutoff of Definition~\ref{def:D13}, every~$b$ and every $K_{\rm rot}\in\mathfrak{so}(4)$,
\begin{equation*}
\mathcal P_b^T=\partial_{x_0^b}S(A(x_0),\boldsymbol B(x)),\qquad
\mathcal Y_{K_{\rm rot}}^T=S((\mathsf L_d(K_{\rm rot})A)(x_0),\boldsymbol B(x))\,.
\end{equation*}
\end{theorem}
The proof is in Section~\ref{sec:partII-G}, II.G.3.1--G.3.2. Divide the raw current by~$c_\star$ and average its tensor transform over~$O(4)$ with probability Haar measure; its shell actions are preserved. The table has full covariance,
\begin{equation*}
S_n(R_{d_i}(h)A_i;hx)=S_n(A;x),\qquad
R_4(h)T_{ab}=\sum_{i,j}h_{ia}h_{jb}T_{ij}\quad(h\in O(4))\,.
\end{equation*}
The free stress~$\maxwellstress$ is normalized by the same shell laws in II.G.1. In particular their sign gives
\begin{equation*}
\maxwellstress_{ab}=-\sum_{c,\alpha}:F_{ac}^\alpha F_{bc}^\alpha:
+\frac{\delta_{ab}}4\sum_{i,j,\alpha}:F_{ij}^\alpha F_{ij}^\alpha:\,.
\end{equation*}

The final comparison includes any chosen stress insertions. For a free list $u_i\in\maxwellspace{d_i}$, choose a subset~$I$ of slots and component pairs $(a_i,b_i)$ on~$I$, and put
\begin{equation*}
A_i^0=\begin{cases}\maxwellstress_{a_ib_i},&i\in I,\\u_i,&i\notin I.\end{cases}
\end{equation*}
The normalized tensor then has the free limit together with all fixed positional derivatives.
\begin{proposition}[Differentiated short-distance comparison]\label{prop:SC}
There is one stable scaling sequence such that, for every fixed list and stress-slot choice, every derivative order~$j$ and compact $\mathcal K\subset\Conf{n}$, with $A_{k,i}=r_k^4T_{a_ib_i}$ on~$I$ and $A_{k,i}=\scaleframe{k}{d_i}u_i$ elsewhere,
\begin{equation*}
D_x^jS_n(A_k;r_kx)\longrightarrow D_x^j\maxwellcorr_n(A^0;x)
\quad\text{uniformly on }\mathcal K\,.
\end{equation*}
\end{proposition}
The proof is in Section~\ref{sec:partII-G}, II.G.3.3. Write $\partial_a^{\rm Max}$ for free field differentiation and $\mathsf L_d^{\rm Max}$ for the generators of the natural rotation action on Wick polynomials. Along this same sequence, for every fixed $d,a,K_{\rm rot}$, the transported infinitesimal actions satisfy the operator-norm limits
\begin{align*}
r_k(\scaleframe{k}{d+1})^{-1}\partial_a\scaleframe{k}{d}&\longrightarrow\partial_a^{\rm Max},\\
(\scaleframe{k}{d})^{-1}\mathsf L_d(K_{\rm rot})\scaleframe{k}{d}&\longrightarrow\mathsf L_d^{\rm Max}(K_{\rm rot})\,.
\end{align*}
In the correlation limit $D_x^j$ differentiates the whole function of~$x$, including the rescaled positions~$r_kx$.

